\documentclass[12pt]{article}
\usepackage{amsmath}
\usepackage{graphicx,psfrag,epsf}
\usepackage{enumerate}
\usepackage{natbib}
\usepackage{url} % not crucial - just used below for the URL 

% Additional packages (kept from original manuscript)
\usepackage{algorithm}
\usepackage{algpseudocode}
\usepackage{tikz}
\usepackage{booktabs}
\usepackage{threeparttable}
\usepackage{tabularx}
\usepackage{array}
\usepackage{multirow}
\usepackage{float}
\usepackage{placeins}
\usepackage[font=small,skip=0pt]{caption}
\usepackage{xcolor}
\usepackage{arydshln}
%\pdfminorversion=4
% NOTE: To produce blinded version, replace "0" with "1" below.
\newcommand{\blind}{0}
\newcommand{\galit}[1]{\textcolor{red}{#1}}
\newcommand{\matt}[1]{\textcolor{blue}{#1}}

% DON'T change margins - should be 1 inch all around.
\addtolength{\oddsidemargin}{-.5in}%
\addtolength{\evensidemargin}{-1in}%
\addtolength{\textwidth}{1in}%
\addtolength{\textheight}{1.7in}%
\addtolength{\topmargin}{-1in}%

\newtheorem{principle}{Principle}
\newenvironment{scenario}[1]{%
  \par\vspace{1ex}\noindent\textbf{#1}\par\vspace{0.5ex}}{\par}

\newtheorem{lemma}{Lemma}
\newtheorem{definition}{Definition}
\newtheorem{remark}{Remark}

\let\oldtable\table
\let\endoldtable\endtable
\renewenvironment{table}
  {\oldtable\spacingset{1}}
  {\endoldtable}
%%%%%%%%%%%%%%%%%%%%%%%%%%%%%%%%%%%%%%%%%%%%%%%%%%%%%%%%%%%%%%%%%%%%%%%%%%%%%%

\begin{document}
%\Supplemental Material %% I don't think this is legal
\clearpage
\bigskip
\begin{center}
{\LARGE \bf Supplemental Material for: \\ When Will Your Model Fail?\\TreeAlert: A Method for Predicting Extreme Forecasting Failures}
\end{center}
\FloatBarrier
\subsection{Example: Probabilistic Forecasting of Heart-Rate}
\label{app:probabilistic_treealert}

To illustrate TreeAlert using a probabilistic forecast, we apply the Winkler score to central 80\% prediction intervals generated from the ARIMA model used in the heart‑rate example\footnote{For the electricity demand example, only point forecasts were made available % because its forecast data is obtained from a 
by the governmental operational platform, and the models themselves are not disclosed. We therefore cannot illustrate probabilistic forecasting for this dataset}% for which full predictive distributions are not publicly available.} 
 (see Section 4.2). Specifically, we set %$(1-\alpha)=0.80$, implying 
$\alpha=0.2$. Next we replace the absolute point‑forecast error with the Winkler score $S_t$ as the outcome variable in the regression tree. This allows TreeAlert to identify patterns of systematically poor interval forecast performance rather than large point-forecast errors alone. 

\FloatBarrier
\subsubsection*{Tree Estimation and Rule Extraction}
The regression tree is estimated on Winkler scores from the training data using the same predictors from Section 4.2 (Minute, Hour, Weekday, Steps, and Intensity). When tuning TreeAlert under the Winkler score objective using grid search, we obtained the same hyperparameter configuration as in the original point‑forecast %heart‑rate 
example. The three terminal nodes with the highest mean Winkler scores are retained and translated into interpretable rules (Table~\ref{tab:winkler_rules}). Each rule corresponds to recurring conditions of poor probabilistic forecast performance arising from overly wide prediction intervals, insufficient coverage, or both.

\begin{table}
\centering
\caption{Interpreted tree rules for terminal nodes with highest mean Winkler scores}
\label{tab:winkler_rules}
\small
\begin{tabular}{c p{12.5cm}}
\hline
\textbf{Avg.~Winkler} & \textbf{Interpreted Rule} \\
\hline
355.56 & The model exhibits high probabilistic forecast loss when Steps\_10 $\geq$ 1139. \\
274.49 & The model exhibits high probabilistic forecast loss when Steps\_10 is 383 to 411, and Hour is 5AM, 2PM - 4PM. \\
231.16 & The model exhibits high probabilistic forecast loss when Steps\_10 is 466 to 1139, Hour is 5AM, 2PM - 4PM, and Minute is 0 through 9, 10 through 19, 20 through 29, and 40 through 49. \\
\hline
\end{tabular}
\end{table}

\subsubsection*{Visualization and Lift Evaluation}

Lift-based evaluation using 20 equal-sized ventiles (Figure~\ref{fig:winkler_lift_test}) yields first-ventile lift of 3.45 in the training period and 3.20 in the test period. These values indicate that the top 5\% of alerted %timestamps ranked 
periods contain more than three times the average Winkler score of the full evaluation period. The close alignment between training and test lift scores suggests that patterns of probabilistic forecast under-performance are generalizable. % to unseen data.

\begin{figure}[H]
  \centering
  \begin{minipage}{0.43\linewidth}
    \centering
    \includegraphics[width=\linewidth]{Figures/winkler_lift_train.png}
  \end{minipage}
  \hfill
  \begin{minipage}{0.43\linewidth}
    \centering
    \includegraphics[width=\linewidth]{Figures/winkler_lift_test.png}
  \end{minipage}
    \caption{Lift chart for heart rate training data (left) and test data (right)}
    \label{fig:winkler_lift_test}
\end{figure}

Figure~\ref{fig:winkler_full} displays the heart rate series with flagged periods corresponding to the retained rules. Flagged timestamps align with periods where observed heart rate values frequently fall near or outside the predictive intervals. Rules 1 and 2 reliably flag periods of elevated Winkler scores in both training and test. Rule 3 flagged periods with low values, suggesting it may have overfit to a false pattern in the training data. 
\begin{figure}[H]
  \centering
  \includegraphics[width=0.9\linewidth]{Figures/winkler_full.png}
  \caption{Heart rate series with flagged periods under Winkler-score TreeAlert}
  \label{fig:winkler_full}
\end{figure}

%\FloatBarrier
\clearpage
\section{Vertical Pruning for Extreme Error Isolation}
\label{append-H}

Figure~\ref{fig:tree} illustrates the ``hollowed out" tree procedure graphically: the shaded region is trimmed to isolate the outer branches corresponding to the most extreme over- and under-forecasting patterns. Retaining only $k$ terminal nodes and their associated rules ensures interpretability by avoiding the complexity of large regression trees that result from high-dimensional features or high-cardinality categorical predictors.  

\begin{figure}[ht]
\centering
\includegraphics[width=.9\textwidth]{Figures/Hollow Tree.png}
\caption{``Hollowed out" regression tree of forecast errors.}
\label{fig:tree}
\end{figure}

\FloatBarrier
%\clearpage

\section{Rule Selection for Over- and Under-forecasting}
\label{append-A}

When domain considerations prioritize predicting both extreme over- and under-forecast failure patterns, TreeAlert can be set to retain both positive and negative mean error rules rather than the most extreme absolute mean errors. 
This section supplements Section 3.4 (Step 4 -- Rule Extraction and Filtering) and Section 3.5 (Step 5 -- Rule Visualization) by illustrating a variant of the Rule Extraction and Filtering step. 

In the electricity demand example, choosing the top-$k$ terminal nodes by absolute mean error resulted in only under-forecast patterns due to their mean magnitudes being most dominant. To reveal over-forecast patterns as well, we partition terminal nodes by the sign of their mean error and then rank separately within each partition. Formally, let $\bar e_j$ denote the mean error for terminal-node $j$. Define $P=\{j:\bar e_j>0\}$, where $P$ is the set of all rules that represent \textit{under-forecasting} patterns (positive mean error), and $N=\{j:\bar e_j<0\}$, where $N$ is the set of all rules that represent \textit{over-forecasting} patterns (negative mean error). We then order $P$ in descending $\bar e_j$ and $N$ in ascending $\bar e_j$, and retain the top $k_1$ nodes from $P$ and the top $k_2$ nodes from $N$ to form the final rule set:
\begin{equation}
\mathcal{R} = \{\, r_j \;|\; j \in P_{k_1} \cup N_{k_2} \,\},
\end{equation}
where $r_j$ denotes the decision rule associated with terminal node $j$, and $P_{k_1}$ and $N_{k_2}$ are the index sets of the top $k_1$ under-forecasting and top $k_2$ over-forecasting nodes, respectively. In this formulation, the total number of retained rules is $k = k_1 + k_2$. 

This alternative filtering method ensures the top-$k$ list is not monopolized by a single error direction and addresses cases where dips are not flagged because their node means are smaller in magnitude than peaks (or vice versa). However, this adjustment comes at the cost of retaining rules that predict less extreme errors.

To illustrate, we revisit the electricity demand demonstration, selecting the three nodes with the highest positive average errors and the three with the lowest negative average errors ($k_1 = k_2 = 3$, yielding $k = 6$) while holding all other parameters constant. These six resulting rules are displayed in Table~\ref{tab:treealert_rules} and visualized in Figure~\ref{fig:3x3}.

\begin{table}
\caption{Tree rules for terminal nodes with top three positive and negative mean errors}
\small
\centering
\label{tab:treealert_rules}
\begin{tabular}{r p{12.5cm}}
\hline
\textbf{Avg.~Error} & \textbf{Rule (Interpreted)} \\
\hline
$4621.59$  & The model under forecasts by an average of 4622 when Day of Week is Sunday through Monday, Hour is 5PM–6PM, and Minute is 0–14. \\
$3757.46$  & The model under forecasts by an average of 3757 when Day of Week is Tuesday through Saturday, Hour is 5PM–7PM, and Minute is 15–44. \\ 
$3626.57$  & The model under forecasts by an average of 3627 when Day of Week is Sunday through Wednesday, Hour is 2AM and 12AM, and Minute is 15–44. \\
$-1398.24$ & The model over forecasts by an average of –1398 when Day of Week is Sunday, Hour is 2PM, and Minute is 0–14 and 30–59. \\
$-606.43$  & The model over forecasts by an average of –606 when Day of Week is Sunday, Hour is 3PM, and Minute is 0–14 and 30–59.\\
$-534.00$  & The model over forecasts by an average of –534 when Day of Week is Sunday, Hour is 1PM and 4PM, and Minute is 0–14.\\
\hline
\end{tabular}
\end{table}

\begin{figure}
  \centering
  \small
  \includegraphics[width=1\linewidth]{Figures/Elec_3x3_both.png}
  \caption{Example of TreeAlert retaining both positive and negative mean error rules}
  \label{fig:3x3}
\end{figure}

\FloatBarrier
%{\spacingset{1.6}
\bibliographystyle{chicago}
\bibliography{references}
%}

% \FloatBarrier
% \section{Configuration and Empirical Performance}
% \label{sec:config-performance}
% \vspace{0.5em}
% \begin{singlespacing}
% \captionof{table}{Summary of TreeAlert configuration and performance across both demonstrations}
% \label{tab:summary}
% \centering
% \setstretch{1}
% \begin{tabular}{lcc}
% \toprule
%  & Electricity Demand & Heart Rate \\
% \midrule
% \textit{Configuration} & & \\
% \quad Temporal predictors & Minute, Hour, Weekday & Minute, Hour, Weekday \\
% \quad Contextual predictors & -- & Steps, Intensity \\
% \quad Train/test split & 75/25 & 85/15 \\
% \quad $n$ & 5 & 9 \\
% \quad $cp$ & 0 & 0 \\
% \quad $k$ & 4 & 3 \\
% \midrule
% \textit{Performance} & & \\
% \quad Training periods flagged & 33 of 2,119 & 55 of 3,240 \\
% \quad Test periods flagged & 16 of 696 & 12 of 572 \\
% \quad $\text{Lift}(g_1)$ -- training & 1.79 & 2.62 \\
% \quad $\text{Lift}(g_1)$ -- test & 1.31 & 2.79 \\
% \bottomrule
% \end{tabular}
% \end{singlespacing}

%\section{TreeAlert's Lift Results for Electricity Forecasts}
%\label{sec:lift_eval_electricity}

%\section{TreeAlert's Lift Results for Heart Rate Forecasts}
%\label{sec:hr_eval_hr}

% %\THEEndNotes
% \begingroup \parindent 0pt \parskip 0.0ex \def\enotesize{\normalsize} \theendnotes \endgroup

% Acknowledgments here
%\ACKNOWLEDGMENT{We extend our sincere gratitude to Ta-Chun Kuan and Carolin Su for their contributions to the development of this research.}

%\bibliographystyle{nonumber}
%%%%%%%%%%%%%%%%%%%%%%%%%%%%%%%%%%%%%%%%%%%%%%%%%%%%%%%%%%%%%%%%%%%%%%
% References
%%%%%%%%%%%%%%%%%%%%%%%%%%%%%%%%%%%%%%%%%%%%%%%%%%%%%%%%%%%%%%%%%%%%%%

\end{document}
